\section{Trabajo futuro}

En esta seccion se detallaran las posibles mejoras y funcionalidades que se pueden implementar para mejorar el proyecto y pueda impactar positivamente para mejorar
sus funcionalidades ya existentes. Algunas de estas mejoras y funcionalidades son:
\vspace{0.25cm}

\subsection{Moderacion de contenido y estrategias de \textit{bloqueo}}
Si bien existe la posibilidad de crear mundos con extrema libertad y lo mismo ocurre con el chat, esto es un problema, 
ya que podria existir la posibilidad de que determinados actores creen mundos que fomenten odio de forma racial,
 etnica o genero. Asi como tambien la creacion de mundos con contenido ilegal o herotico. 
 O utilizar el chat para fomentar el odio o incurrir en practicas de \textit{bullying} o \textit{grooming}.

\vspace{0.25cm}

Para este MVP, se asume que esas cosas no pasan, la realidad es muy adversa y es posible encontrarse con estas prácticas
indebidas. Eso por ello que pensando para un lanzamiento al mercado de nuestro producto, es necesario mitigar estas
practicas a través de moderación y censurar de contenido indebido. Esto puede ser mediante análisis de los contenidos de 
los chats, utilizando herramientas de \textit{Natural Language Processing} (NLP) y moderando los \textit{assets} que se 
suben al juego de forma manual o utilizando herramientas de detección automática con el entrenamiento de modelos.

\vspace{0.25cm}

Para la moderación en chats, tambien pueden ser creados canales de prevención para evitar prácticas de \textit{bullying} 
o \textit{grooming}. Ya se mediante avisos en los chats o interactuando directamente para clausurar el chats a los actores 
que insistan con estas prácticas. Ante la reiterando de acciones malintencionadas en nuestro videojuego. Es necesario utilizar 
estrategias de \textit{bloqueo} de forma temporal o permanente a los jugadores que insistan con estas practicas. Ya se 
mediante de forma manual o con la detección para las diferentes acciones que se mencionaron anteriormente.

\vspace{0.25cm}

\subsection{Herramientas de AI para la generacion de mundos}
La idea es fomentar la creación organica de mundos, pero eso puede volverse una tarea tediosa si el mundo es muy grande 
o muy complejo. Es por eso, que fue considerada una buena idea integrar modelos de AI para que los jugadores 
puedan crear mundos, utilizando estos modelos. Los jugadores escriben prompts y los agentes se 
encargan de crear los mundos en base a los assets disponibles, como asi también crear de forma rapida: 
diálogos de NPCs, ítems dentro del juego o spawns a lo largo del mundo.

\vspace{0.25cm}

De esa forma, se facilitaría la creación de mundos por parte de los creadores. Igualmente hay que destacar que añadir
estas herramientas de IA tiene un consumo energético e hídrico alto, asi que es necesario evaluar esta posibilidad antes
de implementarla.

\vspace{0.25cm}

\subsection{Detección de fraudes}
Dado que la viabilidad del proyecto, viene dada por como se maneja su economia. 
Es importante garantizar que la cantidad de fraudes y estafas sea mínimo para asegurar el beneficio del proyecto y de los 
creadores de mundos. Es por eso que se plantea la posibilidad de añadir modelos preentrenados para facilitar y mitigar la 
detección automáticade fraudes y estafas dentro del videojuego. Detectada estas practicas indebidas, proceder de forma automatica 
y manual con el \textit{bloqueo} de los jugadores que realizan estas practicas y el resarcimiento económico para los jugadores 
y creadores afectados.

\vspace{0.25cm}

\subsection{Migracion del \textit{core-service} hacia microservicios}
Actualmente el Core Service está estructurado internamente como un monolito modular, donde cada módulo 
representa un servicio independiente, pero todos se despliegan como una única unidad. A futuro, se 
contempla migrar hacia una verdadera arquitectura de microservicios, donde cada módulo se despliegue 
de forma independiente, mejorando la disponibilidad y la tolerancia a fallos del sistema.

\vspace{0.25cm}

De esta forma, se puede mejorar cuestiones como la \textit{availability} de nuestra aplicación a través del uso de replicacion
de las diferentes entidades. Facilitaría el desarrollo de nuevos servicios o el mantenimiento de los actuales, sin alterar 
los demás microservicios. Por otro lado, ante una caída de uno de los microservicios, no se alteraria la aplicacion, como 
pasa actualmente con el monolito.

\vspace{0.25cm}

Además de que cada microservicio, necesita una comunicación con los demás servicios, se podrían utilizar herramientas
como: \textit{Kafka} o \textit{RabbitMQ} para persistir los mensajes en caso de caídas y se tiene un mayor auditoria 
sobre los mensajes que pasan entre microservicios. Fomentando la tolerancia a fallos y la seguridad.

\vspace{0.25cm}

\subsection{Migración de los servidores de mundos}
Durante el diseño y la implementación se identificó que Unity y Mirror, si bien son herramientas sólidas 
para juegos multijugador de pequeña escala, no están diseñadas para la optimización y escalabilidad que 
demandan los servidores de un MMO. Por ello, se contempla reescribir el servidor directamente en un 
lenguaje de mayor rendimiento como C++, Go o Rust, lo que permitiría un control más fino sobre los 
recursos del sistema y una mejor adaptación a las exigencias del género.
\vspace{0.25cm}

A su vez, utilizar alguna libreria para la simulación de físicas en los servidores. Al utilizar esto, se podría abstraer 
correctamente las dependencia con estra libreria, permitiendo fácilmente migrar entre las futuras opciones. Como la 
comunicación fue desarrollada con la libreria \textit{protobuf} de \textit{Google}, la migración de la comunicacion entre 
el codigo existe del cliente y el nuevo servidor se realizaría de una forma bastante sencilla y de casi automática.

\vspace{0.25cm}
